% sec1a.tex — c02｜原文 p3–5｜§1 Witten index + §1.1 Supersymmetric harmonic oscillator
% 编号公式 (1.1)–(1.17) 共 17 条；无编号环境。
% [原文 p.3]
\subsection{超对称谐振子}

考虑一个超对称谐振子
\[
L = \frac{1}{2}\dot x^2 - \frac{1}{2}x^2 + \ii\bar\psi\dot\psi - \bar\psi\psi .
\tag{1.1}
\]
该系统的哈密顿量为
\[
H = \dot\psi\,\Pi_\psi + \dot{\bar\psi}\,\Pi_{\bar\psi} + \dot x\,p - L
\tag{1.2}
\]
\[
\phantom{H} = \frac{1}{2}p^2 + \frac{1}{2}x^2 + \bar\psi\psi .
\tag{1.3}
\]
其中 $\Pi_\psi$、$\Pi_{\bar\psi}$ 与 $p$ 分别是 $\psi$、$\bar\psi$ 与 $x$ 的正则共轭量。对系统作量子化，得到产生与湮灭算符
\[
a_B^\dagger = \frac{1}{\sqrt2}(-\ii p + x),\qquad
a_B = \frac{1}{\sqrt2}(\ii p + x),\qquad
a_F^\dagger = \bar\psi,\qquad
a_F = \psi ,
\tag{1.4}
\]
它们满足 $[a_B, a_B^\dagger] = 1$ 与 $\{a_F, a_F^\dagger\} = 1$。用这些算符表出的哈密顿量为
\[
H = (a_B^\dagger a_B + a_F^\dagger a_F) .
\tag{1.5}
\]
不难看出，这个哈密顿量具有超对称性：$Q = a_B^\dagger a_F$，$\bar Q = a_F^\dagger a_B$。这些荷构成的代数为
\[
[Q, H] = [\bar Q, H] = 0 ,\qquad \{Q, \bar Q\} = H .
\tag{1.6}
\]
这是超对称代数的一个普遍特征。除这些连续对称性之外，理论还有一个与我们的讨论相关的离散 $\Z_2$ 对称性，即费米子数 $F$。算符 $F$ 数出态中费米子的数目（模 2）。

谱很容易算出。由于基态满足 $H|0\rangle = 0$，它同时也满足 $Q|0\rangle = 0$。激发态由产生算符作用在真空上得到。
\[
|\chi_B, n\rangle \equiv (a_B^\dagger)^n |0\rangle ,\qquad
|\chi_F, m\rangle \equiv (a_B^\dagger)^m a_F^\dagger |0\rangle .
\tag{1.7}
\]
% [原文 p.4]
第一组态的 $F = 0$，第二组态的 $F = 1$。配分函数也很容易算出：
\[
Z(x) = \Tr x^H = 1 + (x + x^2 + \cdots) + (x + x^2 + \cdots) = \frac{1+x}{1-x} ,\qquad x = \ee^{-\beta} .
\tag{1.8}
\]
让我们详细解释一个可用于计算自由理论配分函数的技巧。如上所述，谱完全由所有产生算符作用在真空上生成。这些产生算符彼此对易，因而独立地作用。方便的作法是在「单个字母」上计算配分函数，也就是在只含单个粒子的态上计算。这称为单字母配分函数，记作 $z(x)$。分别对玻色子与费米子计算它是有用的，二者分别记作 $z_B$ 与 $z_F$。在此情形下，只有一个玻色产生算符和一个费米产生算符。
\[
\begin{aligned}
z(x) &= z_B(x) + z_F(x),\\
z_B(x) &= \Tr_{\text{bosonic letters}} x^H = x,\\
z_F(x) &= \Tr_{\text{fermionic letters}} x^H = x .
\end{aligned}
\tag{1.9}
\]
玻色多粒子态上的配分函数由下式得到
\[
z_B(x) = x \to \frac{1}{1-x} \equiv Z_B(x) ,\qquad
z_F(x) = x \to (1+x) \equiv Z_F(x) .
\tag{1.10}
\]
而完整配分函数为 $Z = Z_B Z_F$。对玻色子，$Z_B$ 对应的是对所有「用任意多个产生算符作用所得态」求和；对费米子，$Z_F$ 对应的是对「没有费米子」与「恰有一个费米子」的态求和。若理论有多种产生算符，单字母配分函数就会有更多项。为了计算完整配分函数，我们希望把 $z_B$ 与 $z_F$ 中的每个单项式都作这一替换，然后取乘积。这由一个称为「齐次指数化」（plethystic exponentiation）的数学操作完成，即 PE。我们定义
\[
PE[f(x_i)] = \exp\left(\sum_{n=1}^{\infty}\frac{1}{n} f(x_i^n)\right),\qquad
\widetilde{PE}[f(x_i)] = \exp\left(-\sum_{n=1}^{\infty}\frac{1}{n}(-1)^n f(x_i^n)\right)
\tag{1.11}
\]
于是
\[
\boxed{\;Z_B(x) = PE[z_B(x)],\qquad
Z_F(x,\mathbf{f}) = \widetilde{PE}[z_F(x,\mathbf{f})],\qquad
Z = Z_B Z_F\;.}
\tag{1.12}
\]

\noindent\textbf{指标}\par
在含相互作用的理论中，我们将无法如此容易地算出配分函数。保持超对称的一般相互作用理论的拉氏量为
\[
L = \frac{1}{2}\dot x^2 - \frac{1}{2}W'(x)^2 + \ii\bar\psi\dot\psi - W''(x)\bar\psi\psi .
\tag{1.13}
\]
% [原文 p.5]
这里函数 $W(x)$ 称为超势。取 $W(x) = x^2/2$ 便回到超对称简谐振子。显然，对任何更高次的多项式 $W(x)$，理论都将是相互作用的。理论的谱、从而配分函数，都难以精确计算。尽管如此，仍存在一个可被精确计算的「受保护量」。
让我们注意到，所有能量 $E \neq 0$ 的态都是二重简并的：一个态是玻色的，另一个是费米子的。这纯粹是超对称的结果。设 $|\psi\rangle$ 是哈密顿量能量为 $E \neq 0$ 的一个本征态，
\[
H(\bar Q|\psi\rangle) = \bar Q H|\psi\rangle = E(\bar Q|\psi\rangle) .
\tag{1.14}
\]
除零能态之外的谱是成对的。这种非零能态的成对现象是只依赖于超对称的稳健现象，即使理论具有保持超对称的相互作用，它也应当成立。当我们打开相互作用时，理论的谱会改变，但由于超对称，各态将始终遵守成对规则。态可以变成零能，但只能成对地变；态也可以离开零能，同样只能成对地离开。这一观察促使 Witten 定义了指标 [10]，
\[
I = \Tr(-1)^F x^H .
\tag{1.15}
\]
看起来这只是在配分函数上作了一个小改动，但差别很大。$(-1)^F$ 因子使非零能处的玻色与费米贡献相互抵消，只留下零能态的贡献。对超对称谐振子，具体地
\[
I = \Tr(-1)^F x^H = 1 + (x + x^2 + \cdots) - (x + x^2 + \cdots) = 1 .
\tag{1.16}
\]
上述论证告诉我们，即使在这个理论中打开保持超对称的相互作用，指标也将始终为 1\footnote{必须确保理论的形变不会导致非归一化的真空态，因为在此类跃迁中，新的零能态可能来自无穷远。关于超对称量子力学语境下的更多讨论，见 [11]。}。换言之，相互作用理论的指标可以通过把相互作用强度设为零来计算，因为指标在光滑形变下不变。这样，前面发展的齐次指数技术就有助于计算指标。
指标只接收 $H = 0$ 的态的贡献；又因为自由理论中能量是可加的，指标也可以从「单字母指标」出发来构造。为计算自由理论的配分函数，我们曾对玻色子与费米子使用两个不同的齐次指数公式。还有一种更简单的办法，从单粒子指标构造多粒子指标：
\[
\boxed{\;I = PE[i(x)],\qquad i(x) = i_B(x) - i_F(x)\;.}
\tag{1.17}
\]
% [Tutor Note: 术语：指标（index）带 $(-1)^F$ 分次、只数零能态；配分函数（partition function）数所有态。这是全篇最易混的一对概念。]
% [Tutor Note: 易错：$(-1)^F$ 中 $F$ 是费米子数算符，玻色态取 +1、费米态取 −1；非零能态成对抵消的正是这一符号。]
% [Tutor Note: 动机：$Q$ 与 $\bar Q$ 把玻色态与费米态互相映射，(1.14) 说明 $E\neq0$ 时 $\bar Q|\psi\rangle$ 必是能量相同的伙伴态。]
% [Tutor Note: 易错：$\Z_2$ 只保留 $F \bmod 2$，因此 $(-1)^F$ 良定义；整数费米子数本身在超对称代数中并不守恒。]
% [Tutor Note: 易错：(1.8) 是普通配分函数、没有 $(-1)^F$；(1.15) 才是指标。自由谐振子两者都可算，但只有后者受保护。]
% [Tutor Note: 术语：逸度 $x = \ee^{-\beta}$，其中 $\beta$ 是 Euclidean 时间圆周长（逆温度）；$x^H$ 即 $\ee^{-\beta H}$。]
% [Tutor Note: 动机：单字母指标 $i(x) = i_B(x) - i_F(x)$ 带相对负号，正是分次的体现；PE 把它组装成多粒子指标。]
% [Tutor Note: 动机：$\widetilde{PE}$ 与 $PE$ 每项差 $(-1)^n$，恰好生成费米子的 $(1+x)$ 展开，从而 (1.10) 成立。]
% [Tutor Note: 文献：Witten 1982 [10] 提出 Witten 指标；脚注 [11] 讨论形变可能导致非归一化真空，需另行排除。]
% [Tutor Note: 易错：(1.12) 中 $\mathbf{f}$ 为原文引入的费米子逸度占位符，但 (1.9)–(1.10) 的 $z_F$ 只依赖 $x$，后文亦未再用；此处照原文保留。]
% [Tutor Note: 术语：superpotential 译「超势」；$W'$ 给出势能、$W''$ 给出费米子质量项，(1.13) 是 $\sN=1$ 量子力学的标准形式。]
% [Tutor Note: 动机：自由理论中能量可加，故多粒子谱是单粒子谱的组合；齐次指数正是这个组合运算的生成函数。]
% [Tutor Note: 语病：原文 p.4 「compute the partition function it so easily」中的 it 疑为衍字，应为「compute the partition function so easily」；照译并登记 ERRATA。]
% [Tutor Note: 语病：原文 p.5 「It is seems like」系动词重复，应为「It seems like」；已登记 ERRATA。]
